\title{Byte Latent Transformer: Patches Scale Better Than Tokens}

\begin{abstract}
We present the Byte Latent Transformer ({\textsc BLT}), a novel architecture for large language models operating directly at the byte level, which achieves—at scale—performance comparable to traditional tokenization-based LLMs, while providing notable gains in inference speed and robustness. {\textsc BLT} organizes raw bytes into variable-length patches that function as its core computational units. These patches are adaptively partitioned according to the predicted entropy of the subsequent byte, ensuring that regions of greater data complexity are allocated more computational resources and model capacity. We conduct the first scaling investigation of byte-level LLMs with controlled floating-point operation (flop) budgets, extending to 8B parameters and 4T training bytes. Our findings validate the practicality of training models on unprocessed bytes with no reliance on a static vocabulary. Efficiency in both training and inference is enhanced by adaptively forming longer patches in sections of predictable data, also leading to better outcomes in reasoning tasks and rare/generalization scenarios. At fixed inference cost, {\textsc BLT} scales much more effectively than tokenization-based approaches, benefitting simultaneously from increased patch length and model capacity.
\end{abstract}

\section{Introduction}
\label{section:intro}

We present the Byte Latent Transformer~(\textbf{BLT}), an architecture that eliminates the need for tokenization by directly processing raw byte sequences. BLT achieves performance on par with traditional tokenization-based models at scale, while also delivering notable gains in computational efficiency and model robustness (\S\ref{sec:robustness}). In contrast, current large language models (llms) are almost completely end-to-end, except for their reliance on tokenization—a heuristic preprocessing procedure that partitions byte streams into a fixed token set. This approach introduces biases in sequence compression and gives rise to limitations such as domain and modality dependence~\citep{dagangetting}, vulnerability to input perturbations~(\S\ref{sec:robustness}), restricted orthographic understanding~\citep{edman2024cute}, and unequal performance across languages~\citep{liang2023xlm,petrov2024language,limisiewicz2024myte}.

Tokenization has historically played a crucial role, as training llms directly on bytes presents significant computational challenges at large scales, primarily due to increased sequence lengths~\citep{xue2022byt5}. Earlier approaches have addressed these issues through more efficient self-attention mechanisms~\citep{el2020characterbert, clark2022canine} or by developing architectures that eliminate attention altogether~\citep{wang2024mambabyte}~(\S\ref{section:related_work}). Nonetheless, such improvements predominantly benefit the training of \textit{small models}. When scaling up, it is the substantial feed-forward network layers—executed across every byte—that account for most of the computational burden in Transformers, rather than the attention modules themselves.

We introduce a dynamic and adaptable approach for organizing bytes into \textit{patches}~(\S\ref{section:patching}), along with a novel model architecture that integrates both byte-level and patch-level data. In contrast to tokenization, {\textsc BLT} does not depend on a predetermined set of patch units. Instead, any collection of bytes can be represented as latent patch vectors through compact, trainable encoder and decoder modules. Our results indicate that this mechanism enables \textit{greater} efficiency in compute distribution compared to traditional tokenization-based architectures.

Traditional tokenization-based language models allocate identical computational resources to each token, prioritizing uniformity over efficiency. This approach relies on compression strategies during token induction that do not always correspond to the actual predictive difficulty associated with different tokens. Our architectural philosophy centers on allowing models to adaptively assign computational effort in response to the complexity of the task at hand. For instance, employing a large transformer is generally unnecessary for predicting the endings of words, as such predictions are typically straightforward and involve less uncertainty than selecting the initial token in a new sentence. The {\textsc BLT} architecture~(\S\ref{section:architecture}) embodies this principle, utilizing three transformer blocks: two compact byte-level \textit{local models} for routine predictions, and a single, larger, global \textit{latent transformer}~(\autoref{fig:arch_high_level}) for more challenging inferences. 

To facilitate dynamic compute allocation, {\textsc BLT} determines how bytes are grouped into patches by analyzing the entropy of subsequent byte predictions, resulting in context-sensitive clusters of bytes that share similar levels of informational content.

We introduce the inaugural study on flop-controlled scaling for byte-level models, investigating configurations with up to 8B parameters trained on 4T bytes, and demonstrate that large-scale, end-to-end training is achievable directly from raw bytes—entirely bypassing the need for fixed-vocabulary tokenization. In our evaluation, {\textsc BLT} achieves flop-normalized training performance\footnote{We determine computational expenditure by tallying the total Floating Point OPerations (flops) performed.} on par with Llama 3, while requiring up to 50\% fewer flops during inference~(\S\ref{section:scaling}).

Moreover, we find that byte-level processing yields substantial benefits for capturing the data's long tail. {\textsc BLT} demonstrates superior resilience to input noise compared to tokenizer-dependent approaches, and exhibits stronger abilities in character-level phenomena, as evidenced by tasks involving orthography, phonological structure, and translation for low-resource languages~(\S\ref{sec:robustness}).

Additionally, using {\textsc BLT} permits joint scaling of both the model’s size and its patch length, without increasing the inference flop budget. Employing longer patch lengths typically reduces computational requirements per inference, allowing more resources to be devoted to expanding the global latent transformer since it is invoked less frequently. In a series of experiments constraining inference flops~(\autoref{fig:fixed_inference_scaling}), we record markedly improved scaling patterns in comparison to traditional tokenization-based models.

To conclude, this work offers several key contributions: 1) We present {\textsc BLT}, a novel byte latent LLM framework that adaptively allocates computational resources, thereby enhancing flop efficiency; 2) Our experiments establish that {\textsc BLT} achieves parity with Llama 3 in terms of training flop consumption up to the 8B parameter level, and enables users to opt for marginal decreases in evaluation performance in exchange for up to 50\% greater flop efficiency; 3) The {\textsc BLT} architecture introduces a new paradigm for scaling LLMs, permitting increases in model size without exceeding a predetermined inference compute budget; 4) We provide empirical evidence that {\textsc BLT} models are more robust to noisy inputs and can capture sub-word features of the input data that escape token-based LLMs. The code for training and inference with {\textsc BLT} is publicly available at~\url{https://github.com/facebookresearch/blt}.

\section{Patching: From Individual Bytes to Groups of Bytes}
\label{section:patching}

Dividing bytes into distinct \textit{patches} enables {\textsc BLT} to flexibly distribute computational resources according to contextual requirements. Figure~\ref{fig:patching_types} illustrates various approaches for partitioning bytes into patches. More precisely, a patching function $f_p$ partitions a byte sequence $\pmb{x} = \{x_i\,|\,i = 1, \ldots, n\}$ of length $n$ into a set of $m < n$ patches, $\pmb{p} = \{p_j\,|\,j = 1, \ldots, m\}$, by assigning each $x_i$ to either 0 or 1, with 1 indicating the initiation of a new patch. In both token-based and patch-based architectures, the overall computation required depends largely on how many steps the core Transformer model must run; in {\textsc BLT}, this correlates to the count of patches generated by the chosen patching function. Therefore, the typical length of a patch, referred to as \textit{patch size}, primarily dictates the computational load for both training and evaluation under any specific patching function~(\S\ref{section:flops}).  
We proceed by describing three patching strategies: dividing data using a fixed byte count per patch~(\S\ref{section:static-patch}), segmenting at whitespace boundaries~(\S\ref{section:white-patch}), and adaptively patching informed by entropy measurements from a lightweight byte-level language model~(\S\ref{section:dyn-patch}). We also address the concept of incremental patching and clarify how the patching process diverges from standard tokenization techniques~(\S\ref{section:bpe-patch}).

\subsection{Strided Patching Every K Bytes}
\label{section:static-patch}

A common and simple approach to grouping bytes is to segment them into uniformly sized patches of length $k$, as illustrated in MegaByte~\citep{yu2023megabyte}. This method, which utilizes a fixed stride, offers practical advantages: it is straightforward to implement during both training and inference stages, enables easy adjustment of average patch size, and allows for clear regulation of computational costs. Nonetheless, there are important limitations to this strategy. Primarily, computation is statically distributed, so resources are not focused where they are most beneficial—for example, transformer steps might be inefficiently used when only whitespace in code is predicted, or may provide insufficient capacity for densely informative byte regions like those containing mathematical content. Additionally, such an approach results in patching that is both inconsistent and insensitive to context, often leading to scenarios where identical byte sequences (such as particular words) are segmented differently.

\subsection{Space Patching}
\label{section:white-patch}

\citet{slagle2024spacebyte} introduces a straightforward yet impactful modification to strided patching, where new patches are initiated following any space-like bytes\footnote{
    Space-like bytes are characterized as any byte that does not correspond to a Latin character, a digit, or a utf-8 continuation byte. Furthermore, at least one byte within each patch must not be space-like.}. These boundaries tend to align with natural linguistic divisions in a variety of languages. The Space patching technique dedicates a latent transformer operation (incurring additional computational cost) to model each word individually. As a result, words are consistently patched across different sequences, and computational resources are focused on challenging prediction points, which commonly arise after spaces. For instance, when generating the initial byte of the response to ``Who composed the Magic Flute?\uline{\hspace{.6em}}'', the model faces a much higher prediction difficulty compared to generating subsequent bytes after the initial ``M'', since the choice of the first character significantly narrows down the possible options, rendering the continuation ``Mozart'' relatively straightforward to infer. Despite these advantages, space patching exhibits limitations in handling certain languages and application domains, and critically, it lacks flexibility in patch size adjustment. The following section presents a novel patching approach informed by the observation that the leading bytes in words pose the greatest predictive challenge, while also naturally enabling adaptable patch sizes.

\subsection{Entropy Patching: Using Next-Byte Entropies from a Small Byte LM}
\label{section:dyn-patch}

In place of using rule-based heuristics like whitespace, we adopt a data-driven strategy to detect instances of heightened uncertainty in next-byte prediction. Our method, termed \textit{entropy patching}, leverages entropy calculations to establish the boundaries of each patch.

To this end, we train a compact byte-level auto-regressive language model using the {\textsc BLT} training dataset and evaluate the next-byte entropy with respect to the language model distribution $p_{e}$ across the byte vocabulary $\mathcal{V}$:

\begin{equation}
H(x_i) = \sum_{v \in \mathcal{V}} p_{e}(x_i=v|\pmb{x}_{<i}) \log p_{e}(x_i=v|\pmb{x}_{<i})
\end{equation}

We investigate two approaches for detecting patch boundaries based on entropies $H(x_i)$. The initial method selects points that exceed a global entropy threshold, as shown in~\autoref{fig:patching}. The alternative method looks for points where entropy values are significantly higher compared to the preceding value. This latter strategy can also be viewed as identifying locations where the generally decreasing entropy pattern within a patch is disrupted.

\begin{align*}
    \text{Global Constraint}\hspace{1cm}H(x_t)& > \theta_g\\
    \text{Approx. Monotonic Constraint}\hspace{1cm}H(x_t)& - H(x_{t-1}) > \theta_r
\end{align*}

Patch boundaries are determined through a streamlined preprocessing procedure carried out at the dataloading stage. This approach contrasts with~\citet{nawrot-etal-2023-efficient}, who employ a classifier trained to detect entropy-based patch boundaries. Within our experimental setup~(\S\ref{section:experiments}), we evaluate and contrast these two techniques for differentiating bytes of low versus high entropy.

\subsection{The Byte-Pair Encoding (BPE) Tokenizer and Incremental Patching}
\label{section:incremental-patching}
\label{section:bpe-patch}

A large number of contemporary llms, such as our baseline Llama 3, employ subword tokenizers based on bpe~\citep{gage1994new, sennrich-etal-2016-neural}. Throughout this work, we define ``tokens'' as groups of bytes selected from a \textit{finite} vocabulary that is established before training. In contrast, we use ``patches'' to describe flexibly grouped byte sequences that are not restricted by a predefined vocabulary. One key distinction between tokens and patches is that, when using tokens, the model does not have explicit access to the original byte-level information.

A significant advancement that {\textsc BLT} introduces over tokenization-based approaches lies in its reconceptualization of the balance between vocabulary size and computational demands. In conventional large language models, enlarging the vocabulary typically results in tokens that encapsulate more bytes on average, thereby reducing the number of processing steps required. However, this also causes an increase in the dimensionality of the output layer, leading to greater computational overhead. This inherent trade-off constrains the extent to which tokenization-based techniques can vary both token length and the cost of inference. As an illustration, Llama 3 expands its average token length from 3.7 bytes to 4.4 bytes, but this enhancement comes at the expense of quadrupling the embedding table's size relative to Llama 2.

During the generation process, {\textsc BLT} must determine, at each step in the byte sequence, whether a patch boundary has been reached, as this dictates whether the Latent Transformer should perform additional computation. This determination must be made in isolation, without reference to parts of the sequence that have not yet been produced. Therefore, the patching mechanism is prohibited from relying on future bytes when deciding how to divide the byte sequence. More precisely, a patching scheme $f_p$ is incrementally consistent if the following holds:
$$f_p(\pmb{x}_{<i}) = f_p(\pmb{x})_{<i}$$
The commonly used bpe algorithm does not qualify as an incremental patching scheme, because it may segment identical prefixes differently depending on what follows in the sequence; consequently, it fails to meet the incremental criterion stated above\footnote{While the introduction of a designated delimiter token for marking patch boundaries can render bpe compatible with incremental patching, this adjustment increases the total number of bytes in the sequence.}.

\section{{\textsc BLT} Architecture}
\label{section:architecture}
The {\textsc BLT} model integrates a large-scale, globally autoregressive language model that processes patch-level representations, complemented by two compact, specialized models: one encodes byte sequences into patch representations, and the other reconstructs bytes from these patch-level embeddings~(\autoref{fig:arch_high_level}).

\subsection{Latent Global Transformer Model}

\textit{The Latent Global Transformer} refers to an autoregressive transformer architecture $\mathcal{G}$ comprised of $l_{\mathcal{G}}$ layers, designed to process a sequence of latent patch inputs, $p_j$, into a corresponding sequence of output patch representations, $o_j$. In this work, indices $j$ and $i$ consistently identify patches and bytes, respectively. The global model is characterized by the application of a block-causal attention mask~\citep{dubey2024llama}, which ensures that attention is confined to patches up to and including the current patch within the same document. Because the majority of floating-point operations during both pre-training and inference are attributed to this component, judiciously selecting when to apply this model enables flexible allocation and regulation of computational resources based on the complexity of specific input and output regions.

\subsection{Local Encoder}
\label{section:local}

\textit{The Local Encoder Model}, represented as $\mathcal{E}$, constitutes a streamlined transformer-based architecture comprising $l_{\mathcal{E}}$ layers, where $l_{\mathcal{E}}$ is substantially smaller than $l_{\mathcal{G}}$. Its principal function is to transform a sequence of input bytes $b_i$ into meaningful patch-level representations $p_j$ in an efficient manner. Notably, this model diverges from standard transformer designs by incorporating a cross-attention mechanism subsequent to every transformer block. This cross-attention layer serves to aggregate the individual byte representations into higher-level patch representations, as illustrated in~(\autoref{fig:crossattn}).

Initially, the byte sequence $b_i$ undergoes embedding via a $\mathbb{R}^{256 \times h_{\mathcal{E}}}$ matrix, resulting in embeddings $x_i$. These initial representations may be further enriched with supplemental hash-embeddings~(\S\ref{sec:hash-emb}), providing the option for additional contextualization.

Next, alternating layers of transformers and cross-attention~(\S\ref{sec:enc-cross-attn}) progressively convert these embeddings into patch representations $p_i$. These patch representations are subsequently fed as input to the global transformer, $\mathcal{G}$, for higher-level processing.

Within the transformer layers, a \textit{local block causal} attention mask is applied: each byte is permitted to attend to a fixed set of $w_{\mathcal{E}}$ immediately preceding bytes. This attention can span across dynamic patch boundaries but is restricted from crossing document boundaries. The subsections below elaborate on both the embedding approach and the design of the cross-attention block.

\subsubsection{Encoder Hash n-gram Embeddings}
\label{sec:ngram-emb}
\label{sec:hash-emb}
An essential factor in building strong and informative representations at each position $i$ involves embedding contextual information from prior bytes. In {\textsc BLT}, this is addressed by capturing features of the single byte $b_i$ on its own as well as when it occurs within an n-gram sequence of bytes. Specifically, at every step $i$, we generate byte-grams

\begin{align}
    g_{i,n}=\{b_{i-n + 1},\ldots, b_i\}
\end{align}

for each byte position $i$ and $n$ ranging from three to eight.\footnote{
We exclude byte-grams of length $n$ or greater whenever $i < n$. }

Next, we propose hash $n$-gram embeddings, where every byte $n$-gram is assigned, via a hash function, to a specific index in an embedding table $E_{n}^{hash}$ with a predefined capacity for every $n \in\{3,4,5,6,7,8\}$~\citep{bai2010learning}.
This embedding is subsequently summed with the embedding of the corresponding byte, followed by normalization, and then supplied as input to the local encoder component. The final composite embedding is given by

\begin{align}
e_i &= x_i + \sum_{n = 3,...,8}E_{n}^{hash}(\text{Hash}(g_{i,n})) \\
\text{where, Hash}(g_{i,n}) &= \text{RollPolyHash}(g_{i,n}) \% |E_{n}^{hash}|
\end{align}

Each $e_i$ is scaled by dividing by the total number of $n$-gram sizes incremented by one, and the RollPolyHash function, as described in Appendix~\ref{appendix:rollpolyhash}, is applied. Section~\ref{section:ablations} presents an ablation study assessing how varying both the parameter $n$ in $n$-gram hash embeddings and the size of the embedding table influences flop-controlled scaling law patterns. Beyond hash-based $n$-gram embeddings, we also conducted experiments with frequency-based $n$-gram embeddings; further information regarding this investigation can be found in Appendix~\ref{appendix:ngrams}.

\subsubsection{Encoder Multi-Headed Cross-Attention}
\label{sec:enc-cross-attn}
Our implementation is largely inspired by the input cross-attention component utilized in the Perceiver architecture~\citep{jaegle2021perceiver}, with a key distinction being that our latent variables are directly linked to individual patch representations rather than employing a predefined, static collection of latents~(\autoref{fig:crossattn}). Each latent is exclusively associated with the bytes forming its respective patch. For each patch $p_j$, a corresponding query vector is constructed by aggregating the byte-level representations of $p_j$ (typically via pooling), which is then processed through a linear transformation, $\mathcal{E}_{C} \in \mathbb{R}^{h_{\mathcal{E}} \times (h_{\mathcal{E}} \times U_{\mathcal{E}})}$, with $U_{\mathcal{E}}$ representing the count of cross-attention heads within the encoder. More formally, if we represent the byte sequence for patch $p_j$ as $f_{\text{bytes}}(p_j)$, then the procedure proceeds as follows

\begin{align}
P_{0,j} &= \mathcal{E}_{C}(f_\text{bytes}((p_j)), f~ \text{is a pooling function} \\
P_l &= P_{l-1} + W_o\left(\text{softmax}\left(\frac{QK^T}{\sqrt{d_k}}\right)V\right) \\
\text{where } Q_j &= W_q(P_{l-1,j}), K_i = W_k(h_{l-1,i}), V_i = W_v(h_{l-1,i}) \\
h_l &= \text{Encoder-Transformer-Layer}_l(h_{l-1})
\end{align}

Here, $P \in \mathbb{R}^{n_p \times h_{\mathcal{G}}}$ denotes the representations of $n_p$ patches that will be input to the global model. These patch representations are constructed by aggregating the corresponding byte embeddings $e_i$ for each patch $p_j$. The matrices $W_q$, $W_k$, $W_v$, and $W_o$ act as linear projections for queries, keys, values, and outputs, respectively, where both the keys and values are derived from projections of byte-level representations $h_i$ from the previous stage (or from $e_i$ when it is the initial layer). A patch-based masking mechanism is employed such that each query $Q_j$ interacts only with the keys and values associated with the specific bytes of patch $j$. Since the model applies multi-headed attention to $Q, K$, and $V$ and patch vectors generally have greater dimensionality ($h_{\mathcal{G}}$) than $h_{\mathcal{E}}$, the patch matrix $P_l$ is maintained with multiple attention heads, each of dimension $h_{\mathcal{E}}$ for the cross-attention step; the outputs are subsequently concatenated to reconstruct the full $h_{\mathcal{G}}$-dimensional patch representations. Furthermore, pre-LayerNorm is applied to queries, keys, and values, and no positional encodings are incorporated in this cross-attention layer. A residual connection is also applied, encompassing the cross-attention computation.

\subsection{Local Decoder} 

The local decoder $\mathcal{D}$, paralleling the design of the local encoder, employs a compact transformer architecture with $l_{\mathcal{D}}$ layers, where $l_{\mathcal{D}}$ is much smaller than $l_{\mathcal{G}}$. Its role is to map the sequence of global patch representations $o_j$ back into the original sequence of raw bytes $y_i$. By leveraging the hidden states generated by the local encoder for the given byte sequence, the local decoder predicts each byte in turn, conditioning each prediction on previously decoded bytes. Its architecture consists of $l_{\mathcal{D}}$ layers arranged in alternating fashion between cross-attention modules and standard transformer layers. In this configuration, the cross-attention layer precedes the transformer block, transforming the patch-level features into byte-level representations, upon which the subsequent transformer layer operates.

\subsubsection{Decoder Multi-headed Cross-Attention} 

Within the decoder's cross-attention mechanism, the assignments of queries and key/value pairs are reversed; the byte-representations serve as queries, while the patch representations function as keys and values. The starting point for the cross-attention's byte-representations is given by the byte embeddings produced by the final encoder layer, specifically $h_{l_{\mathcal{E}}}$. The byte-representations for each following decoder layer $l$, denoted as $d_{l,i}$, are calculated in the following manner:

\begin{align}
D_0 &= h_{l_{\mathcal{E}}} \\
B_l &= D_{l-1} + W_o\left(\text{softmax}\left(\frac{QK^T}{\sqrt{d_k}}\right)V\right), \\
  &\text{where } Q_i = W_q(d_{l-1,i}), K_i = W_k(\mathcal{D}_C(o_j)), V_i = W_v(\mathcal{D}_C(o_j)) \\
  D_l &= \text{Decoder-Transformer-layer}_l(B_l)
\end{align}

In this context, $W_k$ and $W_v$ refer to the projection matrices for keys and values, which apply linear transformations and use the splitting operation $\mathcal{D}_C$ to the output patch embeddings $o_j$ from the global model. The matrix $W_q$ serves as the query projection, operating on byte-level inputs $d_{l-1}$ from the prior layer in the decoder transformer (or on $h_{l_{\mathcal{E}}}$ at the first decoder layer). The final output is passed through an output projection $W_o$, resulting in $B \in \mathbb{R}^{h_{\mathcal{D}} \times n_b}$, where $n_b$ stands for the total number of output bytes. To compute the subsequent decoder states $D_l$, we feed the output of the cross-attention block $B$ into a decoder transformer layer. Consistent with the strategy used for local encoder cross-attention, our setup incorporates multi-headed attention, applies pre-Layer Normalization, omits positional encodings, and incorporates a residual pathway around the cross-attention component.

\section{Experimental Setup}
\label{section:experiments}
Our experimental protocol is meticulously designed to provide a fair comparison between {\textsc BLT} and models employing tokenization. We pay special attention to ensuring that {\textsc BLT} is not granted any undue benefit that could arise from processing longer sequence lengths.

\subsection{Pre-training Datasets} In this study, all model scales are pre-trained using two primary datasets: 1) The Llama 2 dataset~\citep{touvron2023llama}, which consists of 2 trillion tokens sourced from a wide range of publicly accessible origins, subsequently subjected to rigorous cleaning and filtering procedures to enhance data quality; and 2) {\textsc BLT}-1T: a newly constructed dataset comprising 1 trillion tokens from diverse public domains, which also incorporates a portion of the pre-training corpus introduced by Datacomp-LM~\citep{li2024datacomp}. The first dataset is primarily utilized for scaling law analyses, specifically to identify the optimal token count as outlined in~\cite{dubey2024llama}, guiding the selection of optimal architectural parameters for {\textsc BLT}. The second dataset supports a comprehensive pre-training phase for evaluating performance relative to Llama 3 on various downstream benchmarks. Notably, neither dataset contains any information derived from Meta's products or services. Additionally, for baseline assessments involving tokenizer-dependent models, we adopt the Llama 3 tokenizer with a vocabulary capacity of 128K tokens, which consistently resulted in improved baseline outcomes compared to the Llama 2 tokenizer throughout our trials.

\subsection{Entropy Model} For our experiments, the entropy model is implemented as a byte-level language model, which is trained using the same data distribution as the comprehensive {\textsc BLT} model. By default, the architecture employed is a transformer comprising 100M parameters, 14 layers, and a hidden size of 512, equipped with a sliding window attention mechanism covering 512 bytes. All other hyperparameter choices mirror those utilized in both the local and global transformer configurations. We have systematically evaluated the impact of varying model capacities, receptive fields, and architectural designs as detailed in \autoref{section:ablations}. Notably, when the receptive field is sufficiently limited, the resulting entropy model can be compactly represented as a lookup table for efficient encoding.

\subsection{Entropy Threshold and Equalizing Context Length} When applying entropy-based patching, we determine a patching threshold designed to produce a target average \textit{patch size} over the pretraining dataset mixture. In the case of {\textsc BLT}, unlike standard tokenization, the \textit{patch size} is not fixed and can be set arbitrarily, which substantially influences the model's context window. To ensure fair comparisons and prevent models with larger patches from benefiting from longer contexts, we control the expected number of bytes processed per batch. Specifically, this approach requires shortening the sequence length for models that utilize larger patches. For Llama 2 data, the context window is fixed at 8k bytes, whereas for the {\textsc BLT}-1T dataset, we extend the context to an average of 16k bytes, keeping the batch size consistent at an average of 16M bytes.

Although the mean batch size remains unchanged, employing dynamic patching strategies during data loading leads to variability in the proportion of bytes per patch. To enhance computational efficiency, our implementation of {\textsc BLT} training organizes batches such that each contains an equal count of patches, thereby circumventing the need for padding within the more computationally intensive latent transformer. Throughout training, we pad or, if necessary, truncate the input byte sequences to lengths of 12k and 24k bytes for Llama 2 and {\textsc BLT}-1T datasets, respectively. This practice helps mitigate memory usage spikes that might result from exceptionally large patches in some sequences.

\subsection{Entropy Model Context}
\label{section:entropy-context}
Our empirical observations indicate that entropy patching tends to produce increasingly larger patches in highly structured datasets, such as multiple choice formats (refer to \autoref{fig:mmlu_patching} for a representative MMLU example), which characteristically contain significant repetition. These discrepancies stem from the lower entropy values associated with recurrent content within the entropy model context. Accordingly, in the large-scale experiment of {\textsc BLT}-Entropy utilizing a patch size of 4.5, we opted to periodically refresh the entropy context by inserting new lines, and we employed the approximate monotonicity constraint, as it is less vulnerable to "entropy drift" resulting from alterations in context length. It is important to note that this modification influences only the entropy calculation process; the approach used to determine the entropy threshold remains unchanged.

\subsection{FLOPs Estimation} 
\label{section:flops}

Our methodology closely adheres to the transformer floating point operations (flops) calculation outlined by Chinchilla~\citep{hoffmann2022training}, which incorporates the flops required for the feed-forward components, qkvo projections within the self-attention mechanism, as well as for attention calculation and the final output projection. However, we diverge in that we assume the input embedding layer utilizes an efficient lookup rather than a standard dense matrix multiplication, thereby rendering it effectively a 0-flop process. Consistent with prior literature, we approximate that the flops for the backward propagation step are double those of the forward propagation.

To determine the number of flops \textit{per byte} for {\textsc BLT} models, we aggregate the computational costs associated with the local encoder transformer, the global latent transformer, and the local decoder transformer, including the contributions from the cross attention components present in both the encoder and decoder:

\begin{align}
    \text{FL}_{\text{{\textsc BLT}}} &= \text{Transf. FL}(h_{\mathcal{G}}, l_{\mathcal{G}}, m=n_{ctx}/n_p,V=0) / n_p\\
   &+ \text{Transf. FL}(h_{\mathcal{E}}, l_{\mathcal{E}}, m=w_{\mathcal{E}}, V=0) \\
   &+ \text{Transf. FL}(h_{\mathcal{D}}, l_{\mathcal{D}}, m=w_{\mathcal{D}}, V=256) \\ 
    &+ \text{Cross Attn. FL}(h_{\mathcal{E}}, l_{\mathcal{E}}, m=n_p, r=n_p/k) \times k/n_p \\ 
   &+ \text{Cross Attn. FL}(h_{\mathcal{D}}, l_{\mathcal{D}}, m=k, r=k/n_p)
\end{align}

In this context, $n_{ctx}$ denotes the length of the sequence measured in bytes, while $n_p$ specifies the size of each patch. The parameter $r$ represents the ratio between the number of queries and the number of keys/values. The variable $k$ describes the proportion of the patch dimension relative to the byte dimension, effectively indicating the number of separate local model partitions that are merged to generate a unified model representation (with $k=2$ illustrated in \autoref{fig:crossattn}). The symbol $V$ signifies the vocabulary size involved in the output projection, a step that occurs exclusively within the local decoder. The context length for attention, denoted as $m$, varies depending on whether the module operates over the byte-level or patch-level sequence. Consequently, we adjust the attention-related floating point operation counts to fit the specific context of each component. The precise mathematical formulations used for calculating the floating point operations for both the Transformer and Cross-Attention mechanisms are detailed in Appendix~\ref{section:flops-appendix}.

\subsection{Bits-Per-Byte Estimation} The metric of perplexity is inherently tied to a particular tokenizer, as it assesses the model's uncertainty for each token. In situations where models are compared at the byte and token levels, it is common, following prior work~\citep{xue2022byt5, yu2023megabyte, wang2024mambabyte}, to use Bits-Per-Byte (BPB) instead. BPB serves as a tokenizer-agnostic analog of perplexity. Concretely:

\begin{align}
    \text{BPB}(x)&= \frac{\mathcal{L}_{CE}(\pmb{x})}{\ln(2)\cdot n_{\text{bytes}}}
\end{align}

In this formulation, the cumulative cross-entropy loss across the sequence $\pmb{x}$ quantifies the uncertainty in the data, which is then normalized both by the number of bytes in $\pmb{x}$ and a logarithmic factor.

\subsection{Transformer Architecture Hyperparameters} 

Across all transformer blocks within {\textsc BLT}, encompassing both the local and global modules, our implementation largely adopts the Llama~3~\citep{dubey2024llama} configuration. The feed-forward components employ the SwiGLU activation~\citep{swiglu}, while self-attention layers utilize rotary positional embeddings (RoPE)~\citep{RoPe} parameterized with $\theta=500000$~\citep{xiong2024effective}. For normalization across layers, RMSNorm~\citep{RMSNorm} is applied. We leverage Flash attention~\citep{dao2022flashattention} in self-attention layers wherever the attention masks are of the fixed-standard variety, including \textit{block causal} and \textit{fixed-window block causal} masks, with a fixed window size set to 512 tokens in the latter case. In contrast, due to the dynamic, patch-dependent structure of attention in our cross-attention modules, we resort to Flex Attention\footnote{\url{https://pytorch.org/blog/flexattention}}. This choice allows for efficient fused computations, substantially accelerating the training process.

\subsection{{\textsc BLT}-Specific Hyperparameters} 

To evaluate the performance of {\textsc BLT} models, we design experiments addressing both scaling properties and performance on downstream tasks, using models across various sizes: 400M, 1B, 2B, 4B, and 8B. Detailed hyperparameters regarding model architectures can be found in Appendix~Table~\ref{tab:arch}.

For the initial query construction in the first cross-attention module within the local encoder, we apply max-pooling. All {\textsc BLT} models utilize $500,000$ hashes with a single hash function and n-gram sizes spanning from 3 to 8. We set the learning rate uniformly to $4e-4$ for all model sizes. This shared learning rate for both token-based and {\textsc BLT} variants is justified by a hyperparameter search conducted at 400M and 1B model sizes, comparing learning rates from $1e-3$ to $1e-4$, which determined that the chosen value yields optimal results. For Llama-2 scaling trend experiments, training batch sizes recommended in~\cite{dubey2024llama} (or their equivalent in bytes) are adopted.

Optimization is performed using AdamW~\citep{adamw}, with $\beta_1$ fixed at 0.9, $\beta_2$ at 0.95, and $\epsilon=10^{-8}$. A linear warmup is employed for the first 2000 steps, after which a cosine decay schedule is applied to the learning rate, decaying it to 0. Weight decay is set at 0.1, and global gradient norms are clipped to a maximum of 1.0.

\section{Scaling Trends}
\label{section:scaling}

We provide a comprehensive analysis of the scaling behavior of byte-level models to guide the continued scaling of {\textsc BLT} models. Our investigation is designed to overcome shortcomings in earlier studies on byte-level models through several approaches: (a) examining scaling behavior within compute-optimal training paradigms, (b) training aligned 8B parameter models on substantial datasets (up to 1T tokens or 4T bytes) and assessing their performance on downstream benchmarks, and (c) analyzing scaling patterns under fixed inference cost conditions. Subsequent sections will explore particular benefits associated with modeling sequences at the byte level.

\subsection{Parameter Matched Compute Optimal Scaling Trends}
\label{section:parameter-matched-scaling}
We conduct training of several \textit{compute-optimal} bpe and {\textsc BLT} models on the Llama 2 dataset, covering four model scales from 1B up to 8B parameters. Performance in language modeling is then assessed by plotting the number of training flops against results obtained from a representative sample of the data mixture. For the bpe models, the proportion of model parameters to dataset size follows the optimal compute allocation suggested by Llama~3~\citep{dubey2024llama}. This \textit{compute-optimal} configuration is intended, in theory, to deliver the maximal achievable performance for a specified training compute limit~\citep{hoffmann2022training}, thereby serving as a strong point of comparison for our experiments. Alongside each bpe model, we additionally train a {\textsc BLT} variant, ensuring that the Latent Transformer is identical in scale and structure to its corresponding bpe Transformer, and trained on the identical dataset.

As shown in~\autoref{fig:scaling-compute-opt} (right), {\textsc BLT} models consistently achieve performance on par with, or superior to, their bpe-based counterparts as both model size and flops increase. To our knowledge, this work introduces the first byte-level Transformer architecture—{\textsc BLT}—that demonstrates comparable scaling behavior to BPE-based models when operating in compute-optimal settings. These findings support our hypothesis that the parameter-to-training-compute ratio optimal for bpe also largely holds for {\textsc BLT}, or deviates only minimally.

To align with the bpe scaling trajectory, advancements in model architecture and the implementation of dynamic patching are both essential. In ~\autoref{fig:scaling-compute-opt} (left), we present a comparison between space-patching-driven models and Llama 3. For this comparison, we approximate SpaceByte \citep{slagle2024spacebyte} using {\textsc BLT} space-patching, omitting both n-gram embeddings and cross-attention. While SpaceByte offers a performance gain over Megabyte, it still significantly lags behind Llama 3. Moving to \autoref{fig:scaling-compute-opt} (right), we showcase the benefits gained by integrating both architectural refinements and dynamic patching. Notably, {\textsc BLT} models achieve competitive performance with leading tokenizer-based systems like Llama 3, especially at larger scales.

Furthermore, we identify that the selection of tokenizer significantly influences model performance among tokenizer-based models: models trained with the Llama-3 tokenizer consistently surpass those utilizing the Llama-2 tokenizer, even when trained on identical data.

In summary, the {\textsc BLT} architecture occupies a position between Llama 2 and Llama 3, particularly when much larger patch sizes are employed. While Llama 2 and Llama 3's bpe tokenizers have mean token sizes of 3.7 and 4.4 bytes, respectively, {\textsc BLT} demonstrates comparable scaling behaviors with mean patch sizes of 6 and even 8 bytes. Since inference flops are inversely related to the average patch size, adopting an 8-byte patch size nearly halves the inference flops required. Furthermore, increasing the patch size appears to enhance performance as both the model and dataset grow in size. Notably, {\textsc BLT} with an 8-byte patch size initially underperforms relative to bpe Llama 2 at the 1B parameter scale, but surpasses bpe at the 7B scale. This observation points toward the possibility that such large patch sizes could offer greater benefits at even larger scales and that experimenting with even larger patch sizes may be practical as model capacity and computational resources expand.

\subsection{Beyond Compute Optimal Task Evaluations}
\label{section:evals}
To further investigate scaling dynamics, we extend training of an 8B {\textsc BLT} model beyond the compute optimal threshold using the {\textsc BLT}-1T dataset, which is both more extensive and higher in quality. Performance is then evaluated across a collection of well-established benchmarks in both classification and generative settings. For this purpose, we focus on benchmarks that test common sense reasoning, general world knowledge, and code generation capabilities:
\paragraph{Classification tasks} considered are ARC-Easy (0-shot)~\citep{Eval_ARC}, ARC-Challenge (0-shot)~\citep{Eval_ARC}, HellaSwag (0-shot)~\citep{Eval_hellaswag}, PIQA (0-shot)~\citep{Eval_piqa}, and MMLU (5-shot)~\citep{hendrycks2020measuring}. We utilize a prompt-scoring technique that computes the likelihood assigned to choice tokens, reporting the mean accuracy across tasks.
\paragraph{Coding related generation tasks:}  The pass@1 metric is presented for MBPP (3-shot)~\citep{Eval_MBPP} and HumanEval (0-shot)~\citep{Eval_HumanEval}, providing a measure of the model's competence in generating correct Python code solutions.

In \autoref{tab:evals}, we present a comparison among three models trained on the {\textsc BLT}-1T dataset: a baseline model utilizing the bpe Llama 3 tokenizer,\footnote{We select the Llama 3 tokenizer, which features a 128k vocabulary, given its superior performance over Llama 2's 32k vocabulary.} and two {\textsc BLT} model variants. The first employs a space-patching approach ({\textsc BLT}-Space), while the second adopts an entropy-based patching strategy ({\textsc BLT}-Entropy), with an approximate monotonicity constraint and context resets in the entropy-based model upon encountering new lines, as detailed in~\autoref{section:entropy-context}. All models are trained under an identical computational budget in terms of flops. Notably, for {\textsc BLT}-Entropy, we adjust the entropy threshold at inference time from 0.6 to 0.1, which enhances task accuracy but results in a greater number of inference steps.

The {\textsc BLT}-Entropy model achieves superior results compared to the Llama 3 model on 4 of the 7 evaluated tasks, despite both models being trained with an equivalent amount of data (measured in bytes). This performance gain can likely be attributed to two primary factors: (1) more efficient utilization of training computation enabled by dynamic patching, and (2) explicit modeling at the byte level rather than the token level.

Conversely, {\textsc BLT}-Space generally lags behind the Llama 3 tokenizer across nearly all evaluated tasks, except for one. Nevertheless, its increased average patch size of 6 bytes leads to a marked decrease in inference flops. In contrast, both the bpe and entropy-patching models yield similar average patch sizes of about 4.5 bytes on the combined training dataset. Given an identical training budget, the model employing the larger patch size is able to process 30\% more data compared to the other two approaches, which may cause {\textsc BLT} to deviate further from the point of optimal compute efficiency.

\subsection{Patches Scale Better Than Tokens} The {\textsc BLT} architecture enables concurrent scaling of both the model and the patch size, all while keeping the total training and inference computational requirements unchanged and using a fixed quantity of training data. One distinct advantage of patch-based approaches lies in their ability to increase patch size without limitation—something not feasible for token-based models, which are constrained by their fixed-vocabulary design, as elaborated in Section~\ref{section:bpe-patch}. Employing longer patch sizes reduces computational demands, thereby freeing up resources that can be allocated toward enlarging the global latent transformer, which consequently operates less frequently.

A fixed inference scaling analysis is performed to investigate whether increasing model size, paired with reduced step counts on larger patches, leads to improved performance over smaller models using more steps. Beginning with architectures comprising 400m and 3.6B parameters and employing the Llama 2 tokenizer, we identify flop-equivalent models utilizing the Llama 3 tokenizer as well as {\textsc BLT}-Entropy models characterized by mean patch sizes of 6 and 8 bytes on the training datamix (refer to~\autoref{tab:fixed-inf-params} for specific model configurations). For those models with a patch size of 8, three encoder layers are utilized instead of one. Each model is then trained across multiple training flop constraints.

As depicted in \autoref{fig:fixed_inference_scaling}, {\textsc BLT} architectures demonstrate superior scaling performance compared to tokenization-based approaches across both inference flop categories. Notably, while BPE models tend to outperform initially at lower training compute, their advantage diminishes quickly and they are overtaken by {\textsc BLT} models beyond the compute-optimal point. In practical scenarios, allocating more resources during the pretraining phase—despite its one-time cost—can lead to substantially improved model outcomes for a given inference compute constraint. The 8B model class, including systems like Llama 3.1, exemplifies this strategy, having been exposed to data volumes that are two orders of magnitude greater than those required by compute-optimality at that parameter scale.

The point at which {\textsc BLT} surpasses token-based architectures has moved marginally closer to the compute-optimal regime as model size increases within the higher flop class (shifting from 3x to 2.5x the compute-optimal budget). Furthermore, the model using patch size 8 within this higher flop class demonstrates a more pronounced scaling advantage, overtaking competing models earlier. As outlined in Section~\ref{section:parameter-matched-scaling}, increasing the patch size tends to align {\textsc BLT}'s performance with that of BPE-based models as model scale grows. This phenomenon is partially explained by the shrinking proportion of total computational cost attributed to the byte-level Encoder and Decoder modules, which do not scale as rapidly as the Latent Transformer. For example, when total parameter count is expanded 20-fold (from 400M to 8B), the number of {\textsc BLT}'s local model parameters only approximately doubles. This distinction is significant because increasing patch sizes exclusively impacts the computational requirements of the patch Latent Transformer, leaving byte-level module costs unchanged. As a direct consequence, the {\textsc BLT}-Entropy model with ps=8 increased from 1.6x to 1.7x the parameter count of the Llama 2 model when scaling to the larger model class.

To summarize, our investigation into patch-length scaling reveals that the {\textsc BLT} patch-based model architecture exhibits superior scaling behavior when both the patch size and the overall model size are increased together. Notably, these beneficial scaling characteristics continue and even intensify as the model scale grows.

\section{Byte Modeling Improves Robustness} We further evaluate the robustness of {\textsc BLT} in comparison to token-based models that do not utilize explicit byte-level data, and introduce a method for transforming pretrained token-based models into byte-compatible versions.
\label{sec:robustness}

\subsection{Character-Level Tasks}

The initial impetus behind developing byte-level models stemmed from their resilience to noise at the byte level within input data, along with their capacity to recognize the substructure of tokens—an aspect where traditional tokenizer-based architectures often encounter limitations. To systematically assess these qualities, we conduct supplementary experiments on datasets designed to probe both noise robustness and sensitivity to character-level information, spanning English, various languages, as well as numerals and phonemes. The corresponding findings are summarized in Table~\ref{tab:char_tasks}.

\paragraph{Noisy Data} To assess the resilience of {\textsc BLT} in comparison to tokenizer-based models, we generate noisy variants of the benchmark classification datasets described in Section~\ref{section:evals}. Five separate character-level perturbation methods are utilized to introduce noise: (a) \textit{AntSpeak}: The text is reformatted into all uppercase letters, each separated by spaces. (b) \textit{Drop}: 10\% of the text's characters are randomly deleted. (c) \textit{RandomCase}: Randomly alters the casing of characters such that half are uppercase and half are lowercase. (d) \textit{Repeat}: 20\% of characters are randomly selected and each is duplicated up to four times. (e) \textit{UpperCase}: All letters in the text are converted to uppercase. For evaluation, each noise strategy is separately applied to the prompt, the completion, or both, and we present the mean results. Table \ref{tab:char_tasks} displays performance on the noised HellaSwag~\citep{Eval_hellaswag} benchmark, where {\textsc BLT} consistently demonstrates greater robustness than its tokenizer-based counterparts—exhibiting, on average, an 8-point margin over models trained on identical data and also outperforming the Llama 3.1 model trained with substantially more data.

\paragraph{Phonology - Grapheme-to-Phoneme (G2P)} We evaluate how effectively {\textsc BLT} converts sequences of graphemes—written symbols corresponding to a word—into phonemic representations, which capture their pronunciation. Table \ref{tab:char_tasks} summarizes the outcomes of the G2P task conducted in a 5-shot framework utilizing Phonology Bench~\citep{suvarna2024phonologybench}. The results indicate that {\textsc BLT} achieves superior performance compared to the Llama 3 1T tokenizer-based baseline for this task.

\paragraph{CUTE}
To evaluate character-level capabilities, we apply {\textsc BLT} to the CUTE benchmark~\citep{edman2024cute}, which encompasses a variety of tasks organized into three main groups: compositional understanding, orthographic similarity recognition, and sequence manipulation. This benchmark is particularly demanding for most models utilizing tokenizers, as they tend to know the spellings of their tokens yet exhibit difficulty when required to apply this knowledge to actively manipulate text sequences. As shown in Table~\ref{tab:char_tasks}, {\textsc BLT}-Entropy surpasses both BPE Llama 3 models by a margin exceeding 25 points on this benchmark. Notably, our model excels at character manipulation tasks, attaining a near-perfect score of 99.9\% on both spelling-related challenges. These considerable gains, achieved despite {\textsc BLT} being exposed to 16 times less data than Llama 3.1, suggest that character-level competencies are inherently challenging for BPE-based models to acquire. Figure \ref{fig:cute_examples} presents representative examples where the Llama 3 tokenizer-based model falls short, while our {\textsc BLT} model succeeds. Of the tasks considered, only word deletion and insertion see BPE models performing better; however, while such manipulations may pose difficulties for byte-level models, the performance disparity remains modest, and learning to compose words from characters may be more tractable than decomposing words into character-level information. The evaluation protocol remains consistent across all tasks, employing the original Huggingface prompts. It should also be noted that BPE-based models could potentially improve with more carefully designed prompts.

\paragraph{Low Resource Machine Translation} We assess the performance of {\textsc BLT} on translation tasks involving six major language families and twenty-one low-resource languages, spanning a variety of scripts as provided in the FLORES-101 benchmark~\citep{goyal2022flores}. The SentencePiece BLEU scores are summarized in Table~\ref{tab:flores}. Our findings indicate that {\textsc BLT} consistently surpasses the model equipped with the Llama 3 tokenizer, providing a 2-point overall improvement when translating into English and a 0.5-point gain when translating from English. For well-represented language pairs, {\textsc BLT} performs on par with or slightly exceeds the results of Llama 3. Notably, across many low-resource language pairs, {\textsc BLT} achieves stronger results compared to Llama 3, highlighting the advantages of byte-level modeling in effectively handling long-tail byte sequences.

\subsection{Training {\textsc BLT} using Llama 3}
\label{sec:distillation}
We investigate an approach in which {\textsc BLT} models utilize pre-trained tokenizer-based architectures to improve and accelerate training convergence. Specifically, we initialize the global transformer weights in {\textsc BLT} with parameters from a pre-trained Llama 3.1 model. We then fine-tune the global transformer parameters using a learning rate set to one-tenth of that used for the local encoder and local decoder, training for the number of steps optimal for Llama 3. Table~\ref{tab:distill} offers a comparison against a baseline {\textsc BLT}. The results clearly indicate that initializing from Llama 3.1 enables {\textsc BLT} to substantially outperform both the standard Llama 3 and the baseline {\textsc BLT}, with all models trained under equal flop budgets. Furthermore, even when set against the {\textsc BLT}-Entropy model, which was trained on a much more extensive dataset (1T tokens) and reported in Table~\ref{tab:evals}, {\textsc BLT} initialized from Llama 3.1 delivers better performance on the MMLU benchmark. This points to the promise of this methodology for markedly lowering the computational resources needed for training.

This approach may also be interpreted as adapting models that rely on tokenizers into tokenizer-free architectures, thereby transforming a pre-trained LLaMA 3.1 into a {\textsc BLT} variant. For a thorough evaluation, we present results for the original LLaMA 3.1—trained on 15T tokens—in Table \ref{tab:distill}, and compare its outcomes with those achieved by the {\textsc BLT} model derived from LLaMA 3. While our model maintains comparable performance on MMLU and HumanEval, it exhibits more pronounced declines on several other benchmarks. These findings indicate that additional research is necessary to more fully harness the capabilities of the pre-trained backbone and to enhance model outcomes, especially through improved data composition strategies and fine-tuning of hyperparameters.

\section{Ablations and Discussion}
\label{section:ablations}
Here, we present a series of ablation experiments to support the architectural design decisions behind {\textsc BLT}, as well as the chosen patching approach and selected hyper-parameters for the {\textsc BLT} model with 8B parameters trained on the {\textsc BLT}-1T dataset.

\paragraph{Entropy Model Hyper-parameters} We investigate how the scale of the entropy model and the length of the context window influence scaling behavior. For this, we train a suite of byte-level entropy transformer models, with parameter counts ranging from 1 million up to 100 million, and experiment with context window lengths between 64 and 512. Figure \autoref{fig:entropy_ablation} displays the relationships between bits per byte (bpb) and training flop, depicted as scaling law curves. These plots are generated using our $400m$ and $1b$ {\textsc BLT} models, both trained on the Llama-2 dataset. The findings reveal that larger entropy models and longer context windows both improve scaling, although the marginal benefit diminishes for model sizes beyond 50 million parameters.

\paragraph{Types of Patching}

We conduct ablation studies on the four patching strategies described in Section~\ref{section:patching}, specifically: 1) Strided Patching using strides of 4 and 6, 2) Whitespace-based patching, 3) BPE Tokenizer patching leveraging the Llama 3 tokenizer, and 4) Entropy-based patching informed by a lightweight byte-level language model.

Although dynamic patching shortens the sequences' effective length, we adjust for sequence length to ensure that the contextual window remains comparable across all patching methods. To eliminate potential confounds associated with varying context sizes, every model is exposed to an equivalent expected number of bytes per sequence during both training and inference. The outcomes of these controlled experiments are summarized in Figure~\ref{fig:scaling-compute-opt}. Among the approaches evaluated, all alternative patching methods exceed the performance of static patching, with space patching demonstrating performance nearly matching that of dynamic entropy-based patching.

\autoref{tab:evals_ablation} reports the benchmark results for various {\textsc BLT} model variants, including those utilizing tokenizers, space patching, and entropy-based patching, all trained on the Llama 2 dataset for the optimal training duration~\citep{dubey2024llama}. While space patching constitutes a more straightforward approach, as it bypasses the need to execute an entropy model dynamically during training, our results demonstrate that the improvements achieved with entropy-based patching on scaling trends (see Section~\ref{section:scaling}) persist and extend to downstream benchmark performance as well.\footnote{Space patching evaluations derive from prior runs lacking cross-attention; however, similar outcome patterns emerge in settings with cross-attention.}

\paragraph{Cross-Attention} 
In~\autoref{tab:crossattn_ablations_1b}, we perform an ablation study on the impact of incorporating cross-attention at different locations within the encoder and decoder of {\textsc BLT}. For the encoder's cross-attention mechanism, we examine three initialization strategies for the queries: (1) employing an identical learned embedding for all global states, (2) utilizing a hash embedding derived from the bytes present in the patch, and (3) applying a pooling operation to the encoder’s hidden representations of the patch bytes at the specified encoder layer.

Our results indicate that incorporating cross-attention into the \textit{decoder} yields the greatest benefit. Implementing cross-attention within the encoder leads to modest gains, though these are observed exclusively when queries are initialized through pooling. Furthermore, our analysis shows that cross-attention offers notable advantages on the Common-Crawl dataset, with the improvements becoming more pronounced as the patch size increases.

\paragraph{n-gram Hash Embeddings} 

We experiment with n-gram hash embedding vocabularies of sizes 0, 100K, 200K, and 400K, and report the outcomes in Table~\ref{tab:ngram_ablations_1b}. The results indicate that hash embeddings are beneficial across all evaluated domains, with the largest improvements seen in Wikipedia and Github, where the difference reached 0.04 bpb compared to a 0.01 bpb gain after 15k steps at the 8B model scale. Further, at 8B, reducing the number of hashes from 500K to 300K leads to only a negligible performance change of approximately 0.001 bpb at the same step count. These findings suggest that the use of hash embeddings is essential for enabling {\textsc BLT} to achieve performance comparable to models relying on tokenizers. Nonetheless, after incorporating 300K hashes, the performance benefits begin to plateau. Moreover, the data suggests that the advantages of hash embeddings and cross-attention are mostly additive, since they enhance performance on distinct datasets.

\paragraph{Local Model Hyperparamaters} Table \ref{tab:local_layers} presents an ablation of different configurations concerning the layer count for both the local encoder and decoder. The results indicate that, when using hash n-gram embeddings, {\textsc BLT} achieves strong performance with a very minimal encoder—specifically, a single layer—while benefiting from a more substantial decoder.

\section{Related Work}
\label{section:related_work}

\textbf{Character-Level RNNs:} Character language modeling has maintained its significance since the advent of neural network approaches~\citep{sutskever2011generating,mikolov2012subword,graves2013generating}, mainly due to its innate ability to naturally handle out-of-vocabulary terms, thus obviating the need for explicit back-off strategies. For instance, \cite{kim2016character} propose a model that exclusively ingests character-level input through convolutional and highway layers, which in turn are processed by LSTM-based RNNs. This architecture was shown to achieve parity with then-leading RNN language models on English, while yielding superior outcomes for languages with richer morphology—highlighting a key benefit of character-based large language models. In related work, \cite{kenter2018byte} explored byte-level LSTM architectures for machine comprehension, surpassing word-level models for Turkish and Russian, both noted for their complex morphology. Additionally, \cite{zhang2015character} demonstrated the effectiveness of character-level convolutional models for classification problems, where these models occasionally outperformed their word-based counterparts. Building further, \citet{chung2019hierarchical} utilized hierarchical LSTMs with boundary-detection mechanisms at each tier, revealing underlying text hierarchies and yielding further improvements for character-level language modeling. In contrast, ByteNet introduced by \cite{kalchbrenner2016neural} employed CNN-based character-level layers in lieu of attention mechanisms for the task of machine translation.

\textbf{Character-Level Transformers:} The introduction of transformers leveraging attention mechanisms~\citep{vaswani2017attention} in tandem with subword tokenization strategies~\citep{sennrich-etal-2016-neural} brought notable advancements in neural language modeling and established benchmarks. Yet, the use of word and sub-word representations inherently shapes the inductive bias regarding the abstraction layer at which models function. To bridge the capabilities of transformer architectures with early favorable outcomes in character-level language modeling, \citet{al2019character} implemented very deep transformer models. By incorporating auxiliary loss functions, they managed to train transformer-based models that surpassed previous character-level LSTMs. Despite these improvements, there remained a notable performance gap relative to word-level language models. Additionally, GPT-2~\citep{radfordgpt2} reported that, for extensive datasets such as the 1 Billion Word Benchmark, byte-level language models failed to match the effectiveness of their word-level counterparts.

Although \cite{Choe2019BridgingTG} showed that transformer-based LLMs operating at the byte level can surpass subword-level LLMs with similar parameter sizes, such models incur significantly greater computational cost and require substantially longer training times. In a related vein, \cite{el2020characterbert} introduced CharFormer, a BERT-style model that constructs word representations through convolutional layers over character embeddings, reporting notable gains within the medical domain—albeit with increased computational demands. The CANINE model \cite{clark2022canine}, an encoder-only architecture with 150M parameters, is designed to process character sequences directly. CANINE utilizes a deep stack of transformers akin to our global model, coupled with a combination of local transformers and strided convolutions for input character downsampling, and surpasses token-level encoder models like mBERT in multilingual downstream tasks. ByT5~\citep{xue2022byt5} explored fully byte-level encoder-decoder architectures that eschew patching strategies. While ByT5 achieved greater resilience to noise and demonstrated competitive results with tokenizer-based counterparts using one-quarter of the training data, its lack of patching necessitated attention over all input bytes, resulting in high computational overhead. Modeling at the byte level inherently expands sequence lengths, posing challenges to scaling these models efficiently. More recently, \cite{wang2024mambabyte} utilized the Mamba Architecture~\citep{gu2023mamba}, which is able to sustain a fixed-size memory over extremely long contexts, to develop a byte-level Mamba model—also without patching—and demonstrated that, at a scale of 350M parameters, it outperforms comparable byte-level transformer models in flop-matched settings across several datasets in terms of bits-per-byte.

\textbf{Patching-based approaches:} Utilizing patching methods has been shown to significantly reduce the computational overhead—measured in flops—associated with byte-level LLMs, potentially preserving model performance. A range of studies has demonstrated early positive results, particularly at smaller model sizes and with limited training data. For example, \cite{nawrot-etal-2022-hierarchical} introduce the hourglass transformer, which applies static patching for downsampling and upsampling and achieves superior performance over other byte-level baselines at the 150M parameter scale. Expanding on this work, \cite{nawrot-etal-2023-efficient} propose dynamic patching strategies that involve an end-to-end trained boundary predictor, a boundary predictor guided by selected tokenizers, and an entropy-driven patching model akin to {\textsc BLT}. Their experiments show that these dynamic methods can surpass standard transformers of the period on language modeling benchmarks at a 40M parameter scale with approximately 400M tokens. Additionally, \cite{lester2024training} explore the use of sequences compressed via arithmetic coding—which enables higher compression rates than traditional BPE—and leverage equal-info windows to outperform byte-level baselines in language modeling, though these methods still fall short of subword-based baselines.

Our research is primarily influenced by MegaByte~\citep{yu2023megabyte}, a decoder-only causal LLM that leverages a fixed, static patching technique and concatenates latent representations for the transformation of bytes into patches. On the decoder side, MegaByte employs a local model to convert patches back to bytes. Their results indicate that, with 1B parameters, MegaByte achieves parity with tokenizer-based models on datasets comprising approximately 400B bytes. In our experimental analysis, we include comprehensive ablations of MegaByte and observe that its static patching mechanism underperforms compared to compute-optimized, state-of-the-art tokenizer models in a setting where FLOPs are controlled; furthermore, we illustrate that {\textsc BLT} effectively closes this performance gap. In parallel, \cite{slagle2024spacebyte} report similar findings regarding MegaByte, proposing enhancements such as applying patching to whitespaces and similar byte patterns and incorporating a local encoder. Their evaluations show that this approach yields improvements over tokenizer-based transformers in certain domains—including Github and arXiv—at the 1B parameter scale in compute-constrained scenarios. We present additional experiments involving this model and demonstrate that to further advance the scalability and competitiveness of byte-level models—so they rival top-tier token-based models like Llama 3—further architectural innovations are necessary.

\section{Limitations and Future Work}

In this study, our approach to selecting architectural parameters involves training models for the optimal number of steps, as established for Llama~3~\citep{dubey2024llama}. Nonetheless, it is important to note that these scaling laws were derived specifically for BPE-level transformers, which may result in less-than-ideal ratios of data to parameter size when applied to {\textsc BLT}. Determining scaling laws tailored to {\textsc BLT} is deferred to future research and may reveal even more advantageous scaling patterns for our model. Furthermore, most experiments were performed at parameter counts up to 1B, leaving open the possibility that scaling up to 8B parameters and above could alter the optimal architectural configurations and potentially yield further enhancements in performance at these larger model sizes.

Current transformer libraries and codebases are primarily optimized for the high efficiency of tokenizer-based transformer models. Although our work includes theoretically flop-equivalent experiments and leverages efficient implementations like FlexAttention for layers differing from standard transformer designs, the actual runtime performance of our implementations may still lag behind tokenizer-based models. As such, additional optimizations could further improve wall-clock efficiency.

Although {\textsc BLT} currently relies on a distinct entropy model trained separately for patching, integrating the patching model into an end-to-end learning process represents a promising avenue for subsequent exploration. Section \ref{sec:distillation} describes preliminary experiments that provide early evidence for the feasibility of converting tokenizer-based models—such as Llama 3, which are pretrained on over 10T tokens—into byte-based models, by initializing the global transformer with their parameters and then freezing it. Continued research in this area could yield approaches that preserve the advantages of bytefying and even surpass the performance of these tokenizer-based systems, all without necessitating retraining from the ground up.

\section{Conclusion}
\label{section:conclusion}

This work introduces the Byte Latent Transformer (\textbf{BLT}), an architecture that challenges the typical reliance on predetermined token vocabularies within large language models. By implementing a dynamic and trainable strategy for clustering bytes into variable-sized patches, {\textsc BLT} tailors the allocation of computational effort in response to the intrinsic complexity of the input data, thereby achieving notable gains in computational efficiency and model resilience. Through a comprehensive scaling analysis, we show that {\textsc BLT} attains parity with tokenization-based frameworks such as Llama 3 when scaled to models of up to 8B parameters and exposed to 4T bytes of training data, while enabling up to a 50\% decrease in inference computational cost with only marginal detriments to standard evaluation benchmarks. In addition, {\textsc BLT} introduces a new axis for scaling: it facilitates simultaneous expansion of both model and patch dimensions within a constant inference computational envelope—a particularly beneficial property for realistic, resource-constrained environments. Operating directly on byte-level input, {\textsc BLT} also enhances the system’s capability to address rare or outlier data cases, exhibiting heightened resilience to input noise and more nuanced modeling of intra-word features. Collectively, these findings establish {\textsc BLT} as a viable and compelling substitute for conventional tokenization-driven methodologies, presenting a robust, flexible, and scalable paradigm for future language modeling tasks.

\end{document}